
Suppose that the theorem that the table filling algorithm
produces a minimal digraph is false. That is, there is at
least one pair of states $(a, b)$ such that:

\begin{enumerate}
    \item States $a$ and $b$ are distinguishable, in the sense
    that there is some string $w$ that is in exactly one of
    $L_G(a)$ and $L_G(b)$.
    \item The table filling algorithm does not find $a$ and
    $b$ to be distinguished.
\end{enumerate}

Call this pair a bad pair.

If there are bad pairs, then there
must be some that are the shortest among all these strings that
distinguish bad pairs. Let $(p,q)$ be one such bad pair, and let
$w=a_1 a_2 ... a_n$ be a string as short as any that distinguishes
$p$ from $q$. Then $w$ is in exactly one of $L_G(p)$ and $L_G(q)$.

Consider the case where $w=a_1$. Then the table filling algorithm
would have distinguished $p$ from $q$ during the initialization of
the table in the table filling algorithm.

Otherwise let $r$ be an arbitrary neighbor of $p$, $s$ be an arbitrary
neighbor of $q$, and $w'=a_2 ... a_n$. Consider the following cases:

\begin{itemize}
    \item If $w' \notin L_G(r)$ and $w' \notin L_G(s)$, then $w \notin L_G(p)$
    and $w \notin L_G(q)$. This contradicts the assumption that $w$
    is in exactly one of $L_G(p)$ and $L_G(q)$.
    \item If $w' \in L_G(r)$ and $w' \notin L_G(s)$, then the string $w'$ is
    a distinguishing string shorter than $w$, which contradicts the
    assumption that $w$ is a string as short as any which
    distinguishes a bad pair.
    \item If $w' \notin L_G(r)$ and $w' \in L_G(s)$, then the string $w'$ is
    a distinguishing string shorter than $w$, which contradicts the
    assumption that $w$ is a string as short as any which
    distinguishes a bad pair.
    \item If $w' \in L_G(r)$ and $w' \in L_G(s)$, then $w \in L_G(p)$ and
    $w \in L_G(q)$. This contradicts the assumption that $w$ is in
    exactly one of $L_G(p)$ and $L_G(q)$.
\end{itemize}

Thus $r$ and $s$ cannot be a bad pair, which means that the table
filling algorithm would have found $r$ and $s$ to be distinguished.
If $r$ and $s$ could be distinguished, then the algorithm would
have distinguished $p$ from $q$ in the algorithm's next inductive step.
This contradicts the assuption 

\clearpage

Otherwise let $w'=a_2 ... a_n$ and consider the following subcases:
\begin{itemize}
    \item If for all $r \in \delta(p)$ and $s \in \delta(q)$,
    $w' \notin L_G(r)$ and $w' \notin L_G(s)$, then $w \notin L_G(p)$
    and $w \notin L_G(q)$. This contradicts the assumption that $w$
    is in exactly one of $L_G(p)$ and $L_G(q)$.
    \item If there exists an $r \in \delta(p)$ and $s \in \delta(q)$ such that
    $w' \in L_G(r)$ and $w' \notin L_G(s)$, then the string $w'$ is
    a distinguishing string shorter than $w$, which contradicts the
    assumption that $w$ is a string as short as any which
    distinguishes a bad pair.
    \item If there exists an $r \in \delta(p)$ and $s \in \delta(q)$ such that
    $w' \notin L_G(r)$ and $w' \in L_G(s)$, then the string $w'$ is
    a distinguishing string shorter than $w$, which contradicts the
    assumption that $w$ is a string as short as any which
    distinguishes a bad pair.
    \item If for all $r \in \delta(p)$ and $s \in \delta(q)$, $w' \in L_G(r)$
    and $w' \in L_G(s)$, then $w \in L_G(p)$ and $w \in L_G(q)$. This
    contradicts the assumption that $w$ is in exactly one of $L_G(p)$ and
    $L_G(q)$.
\end{itemize}

Thus,


\subsection{Proof of Construction of Minimal Digraph}

We now show the following:

\begin{theorem}
If states $p$ and $q$ are distinguishable, in the sense
that there is some string $w$ that is in exactly one of
$L_G(p)$ and $L_G(q)$, then the table filling algorithm
will find $p$ and $q$ to be distinguished.
\end{theorem}
\begin{proof}
By induction on the length of $w$. As a base case, consider
$w=a_1$. Then the table filling algorithm would have
distinguished $p$ from $q$ during the initialization of the
table in the algorithm.

For the inductive step, assume that $p$
\end{proof}



Suppose that the theorem that the table filling algorithm
produces a minimal digraph is false. That is, there is at
least one pair of states $(a, b)$ such that:

\begin{enumerate}
    \item States $a$ and $b$ are distinguishable, in the sense
    that there is some string $w$ that is in exactly one of
    $L_G(a)$ and $L_G(b)$.
    \item The table filling algorithm does not find $p$ and
    $q$ to be distinguished.
\end{enumerate}

Call this pair a bad pair.

If there are bad pairs, then there
must be some that are the shortest among all these strings that
distinguish bad pairs. Let $(p,q)$ be one such bad pair, and let
$w=a_1 a_2 ... a_n$ be a string as short as any that distinguishes
$p$ from $q$. Then $w$ is in exactly one of $L_G(p)$ and $L_G(q)$.

Consider the case where $w=a_1$. Then the table filling algorithm
would have distinguished $p$ from $q$ during the initialization of
the table in the table filling algorithm.

Otherwise let $w'=a_2 ... a_n$ and consider the following subcases:
\begin{itemize}
    \item If for all $r \in \delta(p)$ and $s \in \delta(q)$,
    $w' \notin L_G(r)$ and $w' \notin L_G(s)$, then $w \notin L_G(p)$
    and $w \notin L_G(q)$. This contradicts the assumption that $w$
    is in exactly one of $L_G(p)$ and $L_G(q)$.
    \item If there exists an $r \in \delta(p)$ and $s \in \delta(q)$ such that
    $w' \in L_G(r)$ and $w' \notin L_G(s)$, then the string $w'$ is
    a distinguishing string shorter than $w$, which contradicts the
    assumption that $w$ is a string as short as any which
    distinguishes a bad pair.
    \item If there exists an $r \in \delta(p)$ and $s \in \delta(q)$ such that
    $w' \notin L_G(r)$ and $w' \in L_G(s)$, then the string $w'$ is
    a distinguishing string shorter than $w$, which contradicts the
    assumption that $w$ is a string as short as any which
    distinguishes a bad pair.
    \item If for all $r \in \delta(p)$ and $s \in \delta(q)$, $w' \in L_G(r)$
    and $w' \in L_G(s)$, then $w \in L_G(p)$ and $w \in L_G(q)$. This
    contradicts the assumption that $w$ is in exactly one of $L_G(p)$ and
    $L_G(q)$.
\end{itemize}

Thus, 